

\noindent\textbf{УДК 519.872}
\vspace{1em}

\begin{center}
{\Large\bfseries Расчёт многоканальных СМО с отрицательными заявками}
\end{center}
\vspace{0.8em}

\begin{center}
Р.\,С.~Хабаров\textsuperscript{1}, к.\,т.\,н.\\[2pt]
E-mail: xabarov.r@yandex.ru\\[6pt]
В.\,А.~Лохвицкий\textsuperscript{1}, д.\,т.\,н.\\[2pt]
E-mail: lokhv\_va@mail.ru\\[8pt]
{\small\textsuperscript{1}Военно-космическая академия им.\,А.\,Ф.\,Можайского, Санкт-Петербург, Россия}
\end{center}
\vspace{0.5em}

\noindent\textbf{Аннотация.} 
Рассматривается расчёт многоканальных систем массового обслуживания с пуассоновскими потоками положительных и отрицательных заявок. Исследуются удаление одной обслуживаемой заявки и мгновенное очищение системы. Для двухфазного гиперэкспоненциального обслуживания описаны матрицы переходов и вычисление стационарных вероятностей методом Такахаси--Таками. Получены выражения для преобразований Лапласа--Стилтьеса времени ожидания, а для модели с очищением --- также времени пребывания. Для других распределений обслуживания используется согласование моментов; комплексные параметры рассматриваются как алгебраическая аппроксимация. Отдельно исследуются приближённые модели повторного обслуживания с новой и с сохранённой длительностью. Проведены проверки по уравнениям баланса, закону Литтла и имитационному моделированию. Указаны ограничения приближений и условия существования моментов.

\textbf{Ключевые слова:} теория массового обслуживания, отрицательные заявки, катастрофы, повторное обслуживание, метод Такахаси--Таками, преобразование Лапласа--Стилтьеса.

\begin{otherlanguage}{english}
\bigskip
\begin{center}\textbf{Calculation of Multiserver Queues with Negative Customers}\end{center}
\textbf{Abstract.} Multiserver queues with independent Poisson streams of positive and negative customers are considered. Negative arrivals either remove one customer in service or clear the entire system. For two-phase hyperexponential service, transition blocks and the Takahashi--Takami stationary calculation are described. Laplace--Stieltjes transforms of waiting times and, for system clearing, sojourn times are obtained. Other service distributions are treated by moment matching; complex parameters provide an algebraic approximation. Restart with resampling and restart with a fixed service requirement are considered separately using approximate effective-service models. Balance equations, Little's law and replicated simulation are used for verification. The limitations of the approximations and moment-existence conditions are stated.

\textbf{Keywords:} queueing theory, negative customers, catastrophes, service restart, Takahashi--Takami method, Laplace--Stieltjes transform.
\end{otherlanguage}


\section*{1. Введение}
Отрицательные заявки описывают воздействия, которые удаляют работу из системы или прерывают её выполнение. Для сетей с положительными и отрицательными заявками известны стационарные распределения в мультипликативной форме при специальных предположениях о потоках, обслуживании и механизме удаления \cite{gelenbe1991,gelenbe1992}. Обзор таких моделей приведён в \cite{artalejo2000}. Для одноканальной системы с произвольным обслуживанием исследованы отрицательные заявки \cite{harrison1996} и мгновенное очищение системы при катастрофе \cite{jain1996}.

В многоканальном случае существенную роль играет распределение остаточных времён обслуживания. Один из подходов к расчёту состоит в фазовой аппроксимации обслуживания и решении уравнений баланса методом Такахаси--Таками \cite{takahashi1976,ryzhikov2018ru}. В настоящей работе этот подход применяется к удалению одной обслуживаемой заявки и к очищению всей системы. Основное внимание уделяется построению переходов и определению временных характеристик. Повторное обслуживание рассматривается отдельно: формулы для времени завершения одной заявки используются в приближённой модели многоканальной системы. Точность этого приближения оценивается численно.

\section*{2. Постановка задачи}
Система имеет \(n\) одинаковых каналов и неограниченную очередь. Положительные и отрицательные заявки поступают по независимым пуассоновским потокам интенсивностей \(\lambda>0\) и \(\delta\ge0.\) Положительные заявки начинают обслуживание в порядке поступления. Исходные длительности обслуживания \(B\) независимы и одинаково распределены, независимы от входных потоков. Обозначим \(b_k=\mathbb E[B^k]\), \(CV=\sqrt{b_2-b_1^2}/b_1.\)

Рассматриваются два механизма воздействия отрицательной заявки:
\begin{enumerate}[leftmargin=1.5em]
\item \textbf{RCS} (Remove Customer in Service): равновероятно выбирается один занятый канал, обслуживаемая заявка удаляется;
\item \textbf{катастрофа}: одновременно удаляются все заявки, находящиеся в обслуживании и в очереди.
\end{enumerate}
В пустой системе отрицательная заявка не оказывает воздействия. При повторном обслуживании удаление заменяется прерыванием: выбранные заявки немедленно начинают обслуживание заново, сохраняя преимущество перед ожидающими. После прерывания либо независимо выбирается новая длительность, либо повторяется первоначальная длительность данной заявки. В обоих случаях выполненная работа теряется; продолжение с места прерывания не рассматривается.

Время ожидания \(W\) отсчитывается от поступления до начала обслуживания, а при очищении --- до начала обслуживания или удаления из очереди. Время пребывания \(V\) заканчивается завершением обслуживания либо удалением. Каждая положительная заявка учитывается один раз. Обозначим \(w_k=\mathbb E[W^k]\), \(v_k=\mathbb E[V^k]\), а через \(q\) --- вероятность завершить обслуживание. Номинальная нагрузка \(\rho=\lambda b_1/n\) используется для задания параметров и не совпадает с фактической занятостью каналов при удалениях. В примерах \(\rho=0{,}7\); для приближённой модели повторного обслуживания дополнительно проверяется \(\lambda\mathbb E[B_{\mathrm{eff}}]<n.\)

\section*{3. Базовая фазовая модель и метод Такахаси--Таками}
\subsection*{3.1. Представление обслуживания и состояния}
Преобразование Лапласа--Стилтьеса (ПЛС) обслуживания аппроксимируется выражением


\[\beta(s)=\sum_{\ell=1}^{2}y_\ell\frac{\mu_\ell}{\mu_\ell+s},\quad
y_1+y_2=1,\]

\[b_k=k!\sum_{\ell=1}^{2}\frac{y_\ell}{\mu_\ell^k},\quad k=1,2,3.\]


При \(y_\ell\ge0\) и \(\mu_\ell>0\) это распределение \(H_2\): фаза выбирается при начале обслуживания и не меняется до его завершения. Если согласование моментов приводит к комплексным параметрам, те же формулы используются алгебраически. Комплексные коэффициенты не являются вероятностями и интенсивностями переходов цепи Маркова; точность такой аппроксимации требует отдельной проверки. Совпадение первых трёх моментов само по себе не гарантирует точности распределения или его хвоста.

Для вещественной фазовой модели состояние задаётся парой \((i,j)\): \(i\) --- число заявок в системе, \(j\) --- число обслуживаемых заявок в фазе 2. Число занятых каналов равно \(m=\min(i,n)\), поэтому \(0\le j\le m.\) Обозначим стационарную вероятность состояния через \(\pi_{i,j}\), вектор вероятностей уровня через \(\boldsymbol\pi_i\), а \(p_i=\sum_j\pi_{i,j}.\) Расчёт ограничивается уровнями \(0,\ldots,N-1\); \(N\) обозначает число учитываемых уровней.

\subsection*{3.2. Блоки переходов}
Матрицы \(A_i\), \(B_i\), \(C_i\) задают переходы с уровня \(i\) соответственно на уровни \(i+1\), \(i-1\) и внутри уровня. Диагональная матрица \(D_i\) содержит суммарные интенсивности выхода. В базовой системе \(C_i=0.\) При \(i<n\) положительный приход даёт переходы


\[(i,j)\to(i+1,j):\ \lambda y_1,\]

\[(i,j)\to(i+1,j+1):\ \lambda y_2.\]


При \(i\ge n\) поступление в очередь не меняет фазы занятых каналов: \((i,j)\to(i+1,j)\) с интенсивностью \(\lambda.\)

На уровнях \(1\le i\le n\) завершение обслуживания задаёт переходы


\[(i,j)\to(i-1,j):\ (i-j)\mu_1,\]

\[(i,j)\to(i-1,j-1):\ j\mu_2.\]


При \(i>n\) освободившийся канал сразу занимает заявка из очереди, поэтому ненулевые элементы строки \(j\) матрицы \(B_i\) имеют вид


\[B_i(j,j)=(n-j)\mu_1y_1+j\mu_2y_2,\]

\[B_i(j,j+1)=(n-j)\mu_1y_2,\qquad
B_i(j,j-1)=j\mu_2y_1.\]


Элементы с индексами вне диапазона \(0,\ldots,n\) отсутствуют. Диагональ определяется суммой выходящих интенсивностей:


\[D_i(j,j)=\lambda+(m-j)\mu_1+j\mu_2,\quad m=\min(i,n).\]


Например, при \(n\ge2\):


\[A_1=\begin{pmatrix}\lambda y_1&\lambda y_2&0\\0&\lambda y_1&\lambda y_2\end{pmatrix},\]

\[B_2=\begin{pmatrix}2\mu_1&0\\\mu_2&\mu_1\\0&2\mu_2\end{pmatrix},\]




\[D_2=\begin{pmatrix}\lambda+2\mu_1&0&0\\0&\lambda+\mu_1+\mu_2&0\\0&0&\lambda+2\mu_2\end{pmatrix}.\]



\subsection*{3.3. Стационарные вероятности и контроль расчёта}
Внутренние уровни удовлетворяют уравнениям


\[\boldsymbol\pi_{i-1}A_{i-1}+\boldsymbol\pi_i(C_i-D_i)+\boldsymbol\pi_{i+1}B_{i+1}=0,\]

\[\sum_i\boldsymbol\pi_i\mathbf1=1.\]


Итерации Такахаси--Таками определяют нормированные векторы уровней \(\boldsymbol t_i\) и отношения \(x_i=p_{i+1}/p_i.\) Используются матрицы \((D_i-C_i)^{-1}\) и произведения с соседними блоками; после сходимости вероятности восстанавливаются рекуррентно и нормируются \cite{takahashi1976,ryzhikov2018ru}.

Сходимость итераций и достаточность усечения проверяются раздельно. Контролируются изменение итерационных величин, невязка уравнений баланса, нормировка, а для вещественной модели --- неотрицательность вероятностей. Затем расчёт повторяется с увеличенным \(N\) и сравниваются целевые характеристики. Малая вероятность последнего уровня сама по себе недостаточна для контроля высоких моментов. При прямых операциях с плотными блоками затраты на проход по уровням имеют порядок \(O(N(n+1)^3)\), объём хранения --- \(O(N(n+1)^2).\)

Для обычной системы \(M/G/n\) с указанной дисциплиной очереди моменты ожидания связаны с факториальными моментами числа ожидающих заявок:


\[w_k\approx\frac{1}{\lambda^k}\sum_{j=k}^{N-1-n}j(j-1)\cdots(j-k+1)p_{n+j},\quad k\ge1.\]


При бесконечной сумме равенство точное, если соответствующий момент конечен. Для удаления и очищения временные характеристики вычисляются по ПЛС, описанным ниже. Независимую проверку первых моментов даёт закон Литтла:


\[\lambda w_1=\sum_i\max(i-n,0)p_i,\]

\[\lambda v_1=\sum_i i p_i,\]


где учитываются все поступившие заявки, включая удалённые.

\section*{4. Модели отрицательных воздействий}
\subsection*{4.1. Удаление одной обслуживаемой заявки}
При \(i>0\) каждый из \(m=\min(i,n)\) занятых каналов выбирается с вероятностью \(1/m.\) Поэтому блок \(B_i\) получается из базового блока заменой \(\mu_\ell\) на \(\mu_\ell+\delta/m\) в интенсивностях завершения обслуживания. Если очередь непуста, удалённую заявку немедленно заменяет следующая. В частности,


\[B_2=\begin{pmatrix}2\mu_1+\delta&0\\\mu_2+\delta/2&\mu_1+\delta/2\\0&2\mu_2+\delta\end{pmatrix}.\]


Для непустых уровней \(D_i^{\mathrm{RCS}}=D_i^{\mathrm{base}}+\delta I\), а \(C_i=0.\) В пустом состоянии диагональ остаётся равной \(\lambda.\) Переходы показаны на рисунке~XREFX0X. Из баланса потока положительных заявок следует


\[q=1-\frac{\delta}{\lambda}(1-p_0).\]



\subsection*{4.2. Мгновенное очищение системы}
В физической модели катастрофа переводит любое непустое состояние непосредственно в \((0,0)\) с интенсивностью \(\delta.\) Чтобы сохранить переходы только между соседними уровнями, вводятся вспомогательные состояния \(\mathcal{D}_i\), \(i\ge1.\) При катастрофе из обычного состояния уровня \(i\ge2\) выполняется переход в \(\mathcal{D}_{i-1}\), а из уровня 1 --- непосредственно в \((0,0).\) Далее \(\mathcal{D}_i\to \mathcal{D}_{i-1}\) с интенсивностью \(\gamma>0\), где \(\mathcal{D}_0=(0,0).\) Положительные заявки во вспомогательных состояниях не поступают.

Схема переходов приведена на рисунке~XREFX1X. Первое состояние каждого ненулевого уровня --- \(\mathcal{D}_i\); за ним следуют обычные фазовые состояния. При таком порядке


\[B_1=\begin{pmatrix}\gamma\\\mu_1+\delta\\\mu_2+\delta\end{pmatrix},\]

\[C_1=\begin{pmatrix}0&0&0\\0&0&0\\0&0&0\end{pmatrix}.\]


Для \(i\ge2\) первый столбец \(B_i\) содержит \(\gamma\) в строке \(\mathcal{D}_i\) и \(\delta\) во всех обычных строках; оставшаяся часть совпадает с базовым блоком обслуживания. Все \(C_i\) равны нулю. Диагональный элемент вспомогательного состояния равен \(\gamma\), обычного состояния --- базовой интенсивности выхода плюс \(\delta.\)

Вспомогательные состояния служат только для решения уравнений. Пусть \(\widetilde\pi\) --- стационарные веса расширенной цепи, а \(c\) --- их сумма по обычным состояниям, включая \((0,0).\) Физические вероятности определяются как


\[c=\widetilde\pi_{0,0}+\sum_{i\ge1}\sum_j\widetilde\pi_{i,j},\]

\[\pi_{i,j}=\widetilde\pi_{i,j}/c.\]


Это соответствует удалению времени пребывания во вспомогательных состояниях: при наблюдении только обычных состояний вся цепочка после катастрофы становится мгновенным переходом в пустую систему. Поэтому для вещественной фазовой модели нормировка восстанавливает исходный процесс при любом \(\gamma>0.\) Масса вспомогательных состояний не прибавляется к вероятности пустой системы. Время ожидания и пребывания также вычисляется только по физическим состояниям. Для комплексной аппроксимации эта операция применяется к алгебраическим весам.

\noindent FIGBLOCKX0X\par

\noindent PANELCAPX0X0X (a) RCS: удаление одной обслуживаемой заявки.\par

\noindent PANELCAPX0X1X (b) Очищение: переход в цепочку вспомогательных состояний.\par

\noindent MAINCAPX0X Рисунок SEQXfig_1X. Переходы от отрицательных воздействий при \(n=3.\) Обозначения внутри обычных состояний показывают числа каналов в фазах 1 и 2. На схеме б стрелка от группы означает переход из каждого её состояния. Переходы от положительных поступлений и завершений обслуживания не показаны.\par

\section*{5. Расчёт временных характеристик}
\subsection*{5.1. Ожидание при удалении одной заявки}
Поскольку вход пуассоновский и независим от обслуживания, поступающая заявка видит стационарное распределение состояний. Пока помеченная заявка ожидает, все \(n\) каналов заняты. В состоянии с \(j\) каналами в фазе 2 интенсивность освобождения канала равна


\[R_j=(n-j)\mu_1+j\mu_2+\delta,\]

\[a_j(s)=\frac{R_j}{R_j+s}.\]


Пусть \(P\) --- матрица вероятностей изменения фаз после освобождения канала и начала обслуживания очередной заявки. Она получается нормировкой строк \(B_{n+1}^{\mathrm{RCS}}.\) Если перед поступлением помеченной заявки в системе было \(k\ge n\) заявок, до её обслуживания нужны \(k-n+1\) освобождений. Поэтому


\[W^*(s)=\sum_{k=0}^{n-1}p_k+
\sum_{k=n}^{N-1}\boldsymbol\pi_k
\left[\operatorname{diag}(\boldsymbol a(s))P\right]^{k-n}\boldsymbol a(s).\]


Вектор \(\boldsymbol a\) здесь является столбцом. Сначала учитывается время пребывания в текущем состоянии, затем переход; этим определяется порядок матриц. Положительные приходы после помеченной заявки на её ожидание не влияют.

Во время обслуживания интенсивность удаления помеченной заявки равна \(\delta/M(t)\), где \(M(t)\) --- текущее число занятых каналов. Замена этой изменяющейся интенсивности постоянной \(r_m=\delta/m\) даёт приближение


\[S_m^*(s)=\frac{r_m}{s+r_m}+\frac{s}{s+r_m}\beta(s+r_m).\]


В расчёте принимается \(m=k+1\), если заявка поступила при \(k<n\), и \(m=n\), если она ожидала. При обозначении положительной части ПЛС ожидания через \(W_+^*(s)\) получаем:


\[V^*(s)\approx\sum_{k=0}^{n-1}p_k S_{k+1}^*(s)+W_+^*(s)S_n^*(s).\]


Для \(n>1\) это приближение сохраняет число занятых каналов на момент начала обслуживания. Оно не является точным преобразованием времени пребывания. Среднее \(v_1\) можно независимо проверить по закону Литтла; при \(n=1\) интенсивность удаления во время обслуживания постоянна, и указанное ограничение исчезает.

\subsection*{5.2. Ожидание и пребывание при очищении}
Пусть \(Y\sim\mathrm{Exp}(\delta)\) --- время до ближайшей катастрофы после поступления помеченной заявки. Введём вспомогательную траекторию, которая начинается из того же стационарного состояния и имеет те же реализации обслуживания, но не содержит последующих катастроф. Обозначим через \(W_0\) время до начала обслуживания помеченной заявки на этой траектории, а через \(Z_0=W_0+B\) --- время до завершения обслуживания.

До ближайшей катастрофы исходная и вспомогательная траектории совпадают. Поэтому непосредственно из механизма очищения следуют равенства


\[W=\min(W_0,Y),\]

\[V=\min(Z_0,Y).\]


Величина \(Y\) независима от \(W_0\) и \(Z_0.\) ПЛС \(W_0^*\) вычисляется по формуле предыдущего подраздела с физическими стационарными весами модели с катастрофами, но без \(\delta\) в интенсивностях освобождения каналов и в матрице \(P.\) При этом \(Z_0^*(s)=W_0^*(s)\beta(s).\)

Для независимых \(T\ge0\) и \(Y\sim\mathrm{Exp}(\delta)\) интегрирование функции выживания даёт


\[\mathbb E[e^{-s\min(T,Y)}]=\frac{\delta}{s+\delta}+\frac{s}{s+\delta}T^*(s+\delta).\]


Следовательно,


\[W^*(s)=\frac{\delta}{s+\delta}+\frac{s}{s+\delta}W_0^*(s+\delta),\]

\[V^*(s)=\frac{\delta}{s+\delta}+\frac{s}{s+\delta}Z_0^*(s+\delta).\]


Вероятность завершения обслуживания равна \(q=Z_0^*(\delta).\) При \(0<q<1\) условные преобразования имеют вид


\[V_{\text{зав}}^*(s)=Z_0^*(s+\delta)/q,\]

\[V_{\text{уд}}^*(s)=\frac{\delta}{s+\delta}\frac{1-Z_0^*(s+\delta)}{1-q}.\]


Индексы обозначают завершение обслуживания и удаление. Эти соотношения точны для заданной вероятностной модели; фазовая аппроксимация и усечение пространства состояний вносят отдельные погрешности.

\subsection*{5.3. Повторное обслуживание с новой длительностью}
Для одной заявки с независимыми пуассоновскими прерываниями интенсивности \(r\) время завершения \(T\) удовлетворяет уравнению восстановления: если \(B<Y\), то \(T=B\); иначе \(T=Y+T'\), где \(Y\sim\mathrm{Exp}(r)\), а \(T'\) --- независимая копия \(T.\) Отсюда \cite{pal2017}


\[B_{\mathrm{eff}}^*(s)=\frac{(s+r)\beta(s+r)}{s+r\beta(s+r)},\]

\[\mathbb E[B_{\mathrm{eff}}]=\frac{1-\beta(r)}{r\beta(r)}.\]


При \(r=0\) используются исходное распределение и его моменты. Для нескольких каналов замена системы обычной \(M/G/n\) с независимыми длительностями \(B_{\mathrm{eff}}\) является приближением. При одновременных прерываниях \(r=\delta\), но общий поток прерываний создаёт зависимость между временами завершения на разных каналах. При RCS индивидуальная интенсивность зависит от числа занятых каналов.

Для RCS в численном приближении используется эвристическое замыкание


\[r=\delta\,\frac{\mathbb E[M]}{\mathbb E[M^2]},\]

\[M=\min(i,n),\]


где распределение \(i\) берётся из вспомогательной системы \(M/G/n\) с эффективным обслуживанием. Это правило не следует из точного усреднения интенсивности удаления и не даёт гарантии точности. Начальное значение \(r=\delta/n\) уточняется итерациями с усреднением старого и нового значений. Расчёт прекращается при изменении менее \(10^{-9}\); при отсутствии сходимости за 100 шагов, бесконечных моментах или неустойчивости вспомогательной системы результат не выдаётся. В устойчивом режиме без удалений \(q=1.\)

\subsection*{5.4. Повторное обслуживание с сохранённой длительностью}
Если длительность \(B=b\) выбирается один раз и повторяется после каждого прерывания, уравнение восстановления применяется условно по \(b.\) Оно даёт


\[\mathbb E[T\mid B=b]=\frac{e^{rb}-1}{r},\]

\[\mathbb E[T^2\mid B=b]=\frac{2e^{rb}}{r^2}(e^{rb}-1-rb).\]


Поэтому при \(M_B(t)=\mathbb E[e^{tB}]\)


\[\mathbb E[T]=\frac{M_B(r)-1}{r},\]

\[\mathbb E[T^2]=\frac{2}{r^2}\bigl(M_B(2r)-M_B(r)-rM'_B(r)\bigr).\]


Для вещественного \(H_2\) первый момент конечен при \(r<\min_{y_\ell>0}\mu_\ell\), второй --- при \(2r<\min_{y_\ell>0}\mu_\ell.\) Для гамма-распределения с масштабом \(\theta\) соответствующие условия имеют вид \(r\theta<1\) и \(2r\theta<1.\) Подстановка за полюсом рационального выражения для \(M_B\) не определяет существующий момент.

В приближённом расчёте многоканальной системы сначала находятся первые два момента \(T\), затем по ним подбирается гамма-распределение, из которого берётся третий момент для фазовой аппроксимации. Этот дополнительный этап не сохраняет точное распределение \(T.\) Для RCS сохраняется указанное выше приближение индивидуальной интенсивности, для общих прерываний --- приближение независимых эффективных обслуживаний. Условия существования моментов при постоянном \(r\) относятся к этой вспомогательной модели; они не являются общим критерием для RCS с изменяющейся интенсивностью.

\subsection*{5.5. Численное получение моментов}
Моменты ПЛС вычисляются по формуле \(\mathbb E[T^k]=(-1)^k(T^*)^{(k)}(0).\) В расчётах удаления и очищения используются центральные конечные разности с девятью узлами и шагом \(h=10^{-3}/b_1.\) Публикуются первые два момента; их значения сопоставляются с независимыми уравнениями и имитацией. Для повторов с новой длительностью применяется деление степенных рядов ПЛС, позволяющее избежать численного дифференцирования. Более высокие моменты без отдельного контроля точности не используются как подтверждённый результат.

\section*{6. Численная проверка}
\subsection*{6.1. Условия эксперимента}
Во всех примерах \(n=3\), \(\lambda=1\), \(b_1=2{,}1\), \(\rho=0{,}7.\) Рассматриваются \(CV=0{,}8\) и \(CV=1{,}2.\) Для удаления и повторов с новой длительностью \(\delta=0{,}3\), для повторов с сохранённой длительностью \(\delta=0{,}02.\) В последнем случае выбор малой интенсивности обеспечивает конечность необходимых моментов вспомогательных распределений в этих примерах.

При \(CV=0{,}8\) имитация использует гамма-распределение со средним \(b_1\) и дисперсией \(b_1^2CV^2\); фазовый расчёт согласует его первые три момента. При \(CV=1{,}2\) и расчёт, и имитация используют вещественное \(H_2\), согласованное с первыми тремя моментами гамма-распределения. Для эффективного обслуживания в формулах повторов используется ПЛС именно выбранного закона: гамма-распределения в первом случае и \(H_2\) во втором.

Для каждой из 12 комбинаций проведено восемь независимых имитационных прогонов. Первые 20\,000 уходов служат прогревом; далее наблюдаются не менее 200\,000 уходов, включая удаления. Генераторы имеют фиксированные различные начальные состояния. Для каждого прогона вычисляются характеристики, затем среднее по прогонам и приближённый 95\%-й доверительный интервал по распределению Стьюдента с семью степенями свободы. Прогрев уменьшает влияние начального пустого состояния, но не заменяет проверку установившегося режима при иных нагрузках.

Расчёт выполнен при \(N=160\) и допуске итераций \(10^{-12}\), затем повторён при \(N=240.\) Максимальное абсолютное изменение \(w_1\) и \(v_1\) составило менее \(10^{-12}.\) Для вещественной модели с удалением и очищением дополнительно решены уравнения стационарной цепи с непосредственными физическими переходами, без использования программных блоков Такахаси--Таками. В базовом примере расхождение вероятностей было менее \(3\cdot10^{-10}\), а средних ожиданий с законом Литтла --- менее \(10^{-8}.\) Первые два момента ожидания проверены также по независимым уравнениям первого шага. Для очищения проверены среднее пребывание и согласованность условных и безусловных моментов.

\subsection*{6.2. Результаты и границы точности}
В таблицах~XREFX2X--XREFX3X приведены расчётные значения и имитационные оценки с полушириной доверительного интервала. Относительное расхождение определяется единообразно: \(\Delta=100(\text{расчёт}-\text{ИМ})/\text{ИМ}.\) Это расхождение с оценкой, а не известная точная погрешность относительно стационарного значения.

В восьми примерах удаления и повторного обслуживания с новой длительностью абсолютное расхождение средних ожиданий не превышает 2{,}14\%, средних пребываний --- 0{,}56\%. Для удаления и очищения все приведённые расчётные средние попадают в соответствующие доверительные интервалы имитации. При \(CV=0{,}8\) сравнение включает погрешность фазовой аппроксимации гамма-распределения, а при \(CV=1{,}2\) обслуживание в обеих моделях совпадает.

\noindent TABLECAPX Таблица SEQXtab_1X. Средние ожидание и пребывание при $\delta=0{,}3$. ДИ --- полуширина 95\%-го доверительного интервала; «общий» означает одновременный повтор на всех занятых каналах.\par

\setlength{\tabcolsep}{3pt}
\begin{tabular}{lcccccc}
\toprule
 & \multicolumn{3}{c}{$w_1$} & \multicolumn{3}{c}{$v_1$} \\
Вариант & Расчёт & ИМ $\pm$ ДИ & $\Delta$, \% & Расчёт & ИМ $\pm$ ДИ & $\Delta$, \% \\
\midrule
\multicolumn{7}{c}{$CV=0,8$} \\
RCS, удаление & 0,4104 & 0,4162 ± 0,0083 & -1,39 & 2,0851 & 2,0907 ± 0,0107 & -0,27 \\
Очищение & 0,1925 & 0,1932 ± 0,0022 & -0,38 & 1,4830 & 1,4835 ± 0,0028 & -0,03 \\
RCS, новый $B$ & 1,2255 & 1,1999 ± 0,0181 & 2,13 & 3,4135 & 3,3957 ± 0,0203 & 0,52 \\
Общий, новый $B$ & 1,8268 & 1,8522 ± 0,0289 & -1,37 & 4,1527 & 4,1759 ± 0,0298 & -0,55 \\
\multicolumn{7}{c}{$CV=1,2$} \\
RCS, удаление & 0,3977 & 0,3998 ± 0,0086 & -0,53 & 1,9223 & 1,9273 ± 0,0107 & -0,26 \\
Очищение & 0,1646 & 0,1642 ± 0,0024 & 0,24 & 1,3117 & 1,3110 ± 0,0036 & 0,05 \\
RCS, новый $B$ & 1,0426 & 1,0465 ± 0,0271 & -0,37 & 3,0441 & 3,0361 ± 0,0311 & 0,26 \\
Общий, новый $B$ & 0,7596 & 0,7722 ± 0,0119 & -1,63 & 2,6509 & 2,6622 ± 0,0109 & -0,43 \\
\bottomrule
\end{tabular}




\noindent TABLECAPX Таблица SEQXtab_2X. Средние при повторе с сохранённой длительностью, $\delta=0{,}02$.\par

\setlength{\tabcolsep}{3pt}
\begin{tabular}{lcccccc}
\toprule
 & \multicolumn{3}{c}{$w_1$} & \multicolumn{3}{c}{$v_1$} \\
Вариант & Расчёт & ИМ $\pm$ ДИ & $\Delta$, \% & Расчёт & ИМ $\pm$ ДИ & $\Delta$, \% \\
\midrule
\multicolumn{7}{c}{$CV=0,8$} \\
RCS, прежний $B$ & 1,0471 & 1,0323 ± 0,0239 & 1,43 & 3,1751 & 3,1615 ± 0,0265 & 0,43 \\
Общий, прежний $B$ & 1,2247 & 1,2506 ± 0,0242 & -2,07 & 3,3994 & 3,4217 ± 0,0259 & -0,65 \\
\multicolumn{7}{c}{$CV=1,2$} \\
RCS, прежний $B$ & 1,6107 & 1,6074 ± 0,0365 & 0,21 & 3,7524 & 3,7549 ± 0,0382 & -0,06 \\
Общий, прежний $B$ & 2,0975 & 2,2168 ± 0,0786 & -5,38 & 4,3113 & 4,4312 ± 0,0847 & -2,71 \\
\bottomrule
\end{tabular}




\noindent TABLECAPX Таблица SEQXtab_3X. Стандартные отклонения ожидания и пребывания. Интенсивности соответствуют таблицам 1 и 2.\par

\setlength{\tabcolsep}{3pt}
\begin{tabular}{lcccccc}
\toprule
 & \multicolumn{3}{c}{$\sigma_W$} & \multicolumn{3}{c}{$\sigma_V$} \\
Вариант & Расчёт & ИМ $\pm$ ДИ & $\Delta$, \% & Расчёт & ИМ $\pm$ ДИ & $\Delta$, \% \\
\midrule
\multicolumn{7}{c}{$CV=0,8$} \\
RCS, удаление & 0,9131 & 0,9340 ± 0,0230 & -2,24 & 1,7043 & 1,7093 ± 0,0157 & -0,29 \\
Очищение & 0,5523 & 0,5555 ± 0,0079 & -0,58 & 1,2655 & 1,2652 ± 0,0052 & 0,02 \\
RCS, новый $B$ & 1,9985 & 1,9601 ± 0,0489 & 1,96 & 2,6943 & 2,6657 ± 0,0347 & 1,07 \\
Общий, новый $B$ & 2,7441 & 2,7967 ± 0,0779 & -1,88 & 3,3921 & 3,4350 ± 0,0638 & -1,25 \\
RCS, прежний $B$ & 1,7708 & 1,7460 ± 0,0447 & 1,42 & 2,4909 & 2,4742 ± 0,0342 & 0,68 \\
Общий, прежний $B$ & 2,0171 & 2,1149 ± 0,0474 & -4,62 & 2,7530 & 2,8203 ± 0,0399 & -2,39 \\
\multicolumn{7}{c}{$CV=1,2$} \\
RCS, удаление & 1,0268 & 1,0252 ± 0,0268 & 0,16 & 2,0976 & 2,0891 ± 0,0159 & 0,40 \\
Очищение & 0,5699 & 0,5696 ± 0,0086 & 0,05 & 1,4337 & 1,4333 ± 0,0040 & 0,03 \\
RCS, новый $B$ & 1,9772 & 1,9836 ± 0,0444 & -0,32 & 3,0613 & 3,0572 ± 0,0319 & 0,14 \\
Общий, новый $B$ & 1,5502 & 1,5817 ± 0,0274 & -1,99 & 2,6415 & 2,6569 ± 0,0181 & -0,58 \\
RCS, прежний $B$ & 2,7956 & 2,8138 ± 0,0911 & -0,65 & 3,8640 & 3,8838 ± 0,0700 & -0,51 \\
Общий, прежний $B$ & 3,4895 & 3,8867 ± 0,2009 & -10,22 & 4,5586 & 4,8699 ± 0,1703 & -6,39 \\
\bottomrule
\end{tabular}





Для повторов с сохранённой длительностью наибольшее расхождение наблюдается при общих прерываниях и \(CV=1{,}2\): около 5{,}38\% по \(w_1\), 2{,}71\% по \(v_1\), 10{,}22\% по стандартному отклонению ожидания и 6{,}39\% по стандартному отклонению пребывания. Эти значения показывают ограничения сведения к независимым эффективным обслуживаниям и дополнительного согласования моментов. Расхождение нельзя объяснять только случайностью имитационной оценки: ряд расчётных значений находится за пределами приведённых интервалов. Полученные результаты не обосновывают равномерную точность при иных нагрузках, интенсивностях прерываний и формах распределения обслуживания.

\FloatBarrier
\section*{7. Заключение}
Описан расчёт стационарных вероятностей многоканальной фазовой системы с удалением одной обслуживаемой заявки и с мгновенным очищением. Для очищения предложено использовать вспомогательные состояния с последующим исключением времени пребывания в них. Это сохраняет структуру соседних уровней и даёт физическую нормировку вероятностей. Формулы ожидания учитывают время до последовательных освобождений каналов; формулы пребывания при очищении следуют из совпадения траекторий до ближайшей катастрофы.

Для повторного обслуживания рассмотрена приближённая модель эффективной длительности. Формулы для одной заявки с постоянной интенсивностью прерываний отделены от дополнительных приближений многоканальной системы. Численная проверка подтверждает применимость расчёта в рассмотренных примерах, но выявляет более заметные расхождения при общих прерываниях с сохранением длительности. Уточнение зависимости между каналами и интенсивности прерывания помеченной заявки представляет направление дальнейшей работы. Программная реализация основана на библиотеке Most-Queue \cite{mostqueue2026}; параметры, начальные состояния генераторов и результаты прогонов сохранены вместе с материалами расчёта.


